% 原始模板规范：所有学位类型引用均为双重上标格式
\newcommand{\upcite}[1]{\textsuperscript{\textsuperscript{\cite{#1}}}}

\newcommand{\underlineFixlen}[2][3.5cm]{\underline{\makebox[#1][c]{#2}}}

\newlength{\cauc@saved@headsep}
\newcommand{\cauc@abstractspacing}{%
    \ifx\cauc@linespacing\@empty
        % 原始模板规范：所有学位类型摘要均使用baselineskip=20pt(stretch≈1.389)
        1.389%
    \else
        \cauc@linespacing
    \fi
}

\newcount\cauc@abstract@lines

\newenvironment{abstract-zh}{
    \setlength{\cauc@saved@headsep}{\headsep}
    \ifcauc@bachelor
        \begin{center}
            \zihao{3}\fakehei 摘\quad 要
        \end{center}
    \else
        \chapter*{\heiti \texorpdfstring{摘\qquad 要}{摘要}}
        \phantomsection
        \addcontentsline{toc}{chapter}{\texorpdfstring{摘\qquad 要}{摘要}}
    \fi
    \markboth{}{}
    \setlength{\headsep}{1.224cm}
    % 原始模板规范：本科用\begin{center}排版标题，无afterskip，需手动间距；研究生用\chapter*，已有afterskip=30pt
    \ifcauc@bachelor
        \vspace{21pt}%
    \fi
    \zihao{-4}\songti
    \cauc@abstract@lines=0%
    \everypar{\advance\cauc@abstract@lines by \prevgraf}%
    \begin{spacing}{\cauc@abstractspacing}
}{%
    \end{spacing}%
    \advance\cauc@abstract@lines by \prevgraf
    % 中文摘要阈值10行（约300字），依据学位论文规范最低字数要求
    \ifnum\cauc@abstract@lines<10
        \cauc@warning{Chinese abstract appears too short (approximately \the\cauc@abstract@lines\space lines, expected at least 300 characters / 10 lines)}%
    \fi
    \everypar{}%
    \setlength{\headsep}{\cauc@saved@headsep}%
}

\newenvironment{abstract-en}{
    \setlength{\cauc@saved@headsep}{\headsep}
    \ifcauc@bachelor
        \begin{center}
            \zihao{3}\textbf{Abstract}
        \end{center}
    \else
        % 原始模板规范：研究生英文摘要标题使用\textbf（粗体），与中文摘要标题\heiti不同
        \chapter*{\normalfont\zihao{3}\textbf{Abstract}}
        \phantomsection
        \addcontentsline{toc}{chapter}{Abstract}
    \fi
    \markboth{}{}
    \setlength{\headsep}{1.224cm}
    % 原始模板规范：本科需手动间距，研究生\chapter*已有afterskip
    \ifcauc@bachelor
        \vspace{21pt}%
    \fi
    \zihao{-4}
    \cauc@abstract@lines=0%
    \everypar{\advance\cauc@abstract@lines by \prevgraf}%
    \begin{spacing}{\cauc@abstractspacing}
}{%
    \end{spacing}%
    \advance\cauc@abstract@lines by \prevgraf
    % 英文摘要阈值6行（约100词），依据学位论文规范最低字数要求
    \ifnum\cauc@abstract@lines<6
        \cauc@warning{English abstract appears too short (approximately \the\cauc@abstract@lines\space lines, expected at least 100 words / 6 lines)}%
    \fi
    \everypar{}%
    \setlength{\headsep}{\cauc@saved@headsep}%
}

\newcount\cauc@keyword@count
\newcommand{\cauc@check@keywords}[1]{%
    \cauc@keyword@count=1%
    \cauc@count@sep#1；\@nil
    \ifnum\cauc@keyword@count<3
        \cauc@warning{Keywords count (\the\cauc@keyword@count) is less than 3}%
    \fi
}
\def\cauc@count@sep#1；#2\@nil{%
    \ifx\relax#2\relax\else
        \advance\cauc@keyword@count by 1\relax
        \cauc@count@sep#2\@nil
    \fi
}

\newcommand{\keywordszh}[1]{
    % 原始模板规范：本科用\newline\newline，研究生用\newline；\leavevmode确保在垂直模式下也可用
    \ifcauc@bachelor
        \leavevmode\newline\leavevmode\newline\indent{\bfseries\zihao{-4} 关键词：}\zihao{-4}\songti #1
    \else
        \leavevmode\newline
        \noindent{\bfseries\zihao{-4} 关键词：}\zihao{-4}\songti #1
    \fi
    \cauc@check@keywords{#1}
}

\newcommand{\keywordsen}[1]{
    % 原始模板规范：本科用\newline\newline，研究生用\newline；博士英文关键词不使用\songti（默认Times New Roman）
    \ifcauc@bachelor
        \leavevmode\newline\leavevmode\newline\indent{\bfseries\zihao{-4} Key Words: }\zihao{-4}\songti #1
    \else
        \leavevmode\newline
        \noindent{\bfseries\zihao{-4} Key words: }\zihao{-4}%
        \ifcauc@doctor\else\songti\fi #1
    \fi
}

\newenvironment{acknowledgment}{
    \chapter*{\heiti \texorpdfstring{致\qquad 谢}{致谢}}
    \phantomsection
    \addcontentsline{toc}{chapter}{\texorpdfstring{致\qquad 谢}{致谢}}
    \markboth{}{}
    \zihao{-4}\songti
}{}

% 研究生作者简介，原始模板规范要求
\ifcauc@graduate
\newenvironment{authorintro}{
    \chapter*{\heiti 作者简介}
    \phantomsection
    \addcontentsline{toc}{chapter}{作者简介}
    \markboth{}{}
    \zihao{-4}\songti
}{}
\fi

\newif\ifcauc@tabularx
\newcommand{\cauc@if@tabularx}[1]{%
    \cauc@tabularxfalse
    \@tfor\cauc@temp:=#1\do{%
        \expandafter\ifx\cauc@temp X\cauc@tabularxtrue\fi
    }%
}
\NewDocumentCommand{\threelinetable}{m m m m m}{
    \begin{table}[!ht]
        \centering
        \caption{#2}\label{#3}
        \cauc@if@tabularx{#1}%
        \ifcauc@tabularx
            \begin{tabularx}{\textwidth}{#1}
        \else
            \begin{tabular}{#1}
        \fi
        \toprule
        #4
        \midrule
        #5
        \bottomrule
        \ifcauc@tabularx
            \end{tabularx}
        \else
            \end{tabular}
        \fi
    \end{table}
}

\RequirePackage{amsthm}

\theoremstyle{plain}
\newtheorem{theorem}{定理}[chapter]
\newtheorem{lemma}[theorem]{引理}
\newtheorem{corollary}[theorem]{推论}
\newtheorem{proposition}[theorem]{命题}

\theoremstyle{definition}
\newtheorem{definition}[theorem]{定义}
\newtheorem{example}[theorem]{例}

\theoremstyle{remark}
\newtheorem{remark}[theorem]{注}
\renewcommand{\proofname}{证明}

\counterwithin{figure}{chapter}
\counterwithin{table}{chapter}
\counterwithin{equation}{chapter}
% 在 \counterwithin 之后设置编号格式，确保不被覆盖
% 本科：图/表/公式用 - 分隔（如 1-1），研究生：表/公式用 . 分隔（如 1.1）
\renewcommand{\thefigure}{\arabic{chapter}-\arabic{figure}}
\ifcauc@bachelor
\renewcommand{\thetable}{\arabic{chapter}-\arabic{table}}
\renewcommand{\theequation}{\arabic{chapter}-\arabic{equation}}
\else
\renewcommand{\thetable}{\arabic{chapter}.\arabic{table}}
\renewcommand{\theequation}{\arabic{chapter}.\arabic{equation}}
\fi
\ifcauc@graduate
\counterwithin{algocf}{chapter}
\renewcommand{\thealgocf}{\arabic{chapter}.\arabic{algocf}}
\fi

\ifcauc@graduate
\makenomenclature

\renewcommand{\nomname}{\centerline{\heiti\zihao{3}符号说明}}

\renewcommand{\nompreamble}{}
\renewcommand{\nompostamble}{}

\setlength{\nomlabelwidth}{2.5cm}
\setlength{\nomitemsep}{0pt}

\renewcommand{\nomgroup}[1]{%
  \ifx#1A\item[\textbf{\heiti 缩略语}]%
  \else\item[\textbf{\heiti 符号}]%
  \fi
}

\renewcommand{\nomlabel}[1]{\textsf{#1}\hfill}

% 原始模板规范：研究生四级标题使用\paragraph+换行，标题独占一行，正文另起一行
\newcommand{\subsubsubsection}[1]{\paragraph{#1}\mbox{}\\}
\setcounter{secnumdepth}{4}

\renewcommand\thefootnote{\zihao{-5}\ding{\numexpr171+\value{footnote}}}

\bibpunct{[}{]}{,}{n}{,}{,}

\NewDocumentEnvironment{thesislongtable}{m m m m m o}{%
    \setlength{\abovecaptionskip}{0cm}
    \renewcommand\arraystretch{1.2}
    \begin{longtable}{#4}
        \bicaption{#1}{#2}\label{#3}\\
        \toprule
        #5 \\
        \midrule
    \endfirsthead
        \bicaption{#1}{#2}（续）\\
        \toprule
        #5 \\
        \midrule
    \endhead
        \midrule
    \endfoot
        \bottomrule
    \endlastfoot
}{%
    \end{longtable}
    \IfNoValueF{#6}{%
        \linespread{1.0}\setlength{\parindent}{2em}
        \par\smallskip
        #6
    }
    \vspace{-10pt}
}

\fi

\ifcauc@bachelor
\providecommand{\printnomenclature}{%
    \ClassWarning{cauc_thesis}{printnomenclature is only available in graduate mode}%
}
\bibpunct{[}{]}{,}{n}{,}{,}
\fi

\NewDocumentEnvironment{thesistable}{m m m m m o o}{%
    \setlength{\abovecaptionskip}{0cm}
    \begin{table}[!ht]
        \renewcommand\arraystretch{1.2}
        \ifcauc@bachelor
            \caption{#1}
        \else
            \bicaption{#1}{#2}
        \fi
        \label{#3}
        \vspace{-10pt}
        \begin{center}
        \cauc@if@tabularx{#4}%
        \ifcauc@tabularx
            \IfNoValueTF{#7}{%
            \begin{tabularx}{\textwidth}{#4}%
            }{%
            \begin{tabularx}{#7}{#4}%
            }%
        \else
            \begin{tabular}{#4}%
        \fi
            #5
        \ifcauc@tabularx
            \end{tabularx}
        \else
            \end{tabular}
        \fi
        \end{center}
        \vspace{-5pt}
        \IfNoValueF{#6}{%
            \linespread{1.0}\setlength{\parindent}{2em}
            \par\smallskip
            #6
        }
        \vspace{-10pt}
    \end{table}
}{}

\NewDocumentEnvironment{thesisfigure}{m m m O{0.8\textwidth} o}{%
    \setlength{\abovecaptionskip}{0cm}
    \begin{figure}[!ht]
        \ifcauc@bachelor
            \caption{#1}
        \else
            \bicaption{#1}{#2}
        \fi
        \label{#3}
        \vspace{-10pt}
        \begin{center}
        \setlength{\linewidth}{#4}
}{%
        \end{center}
        \vspace{-5pt}
        \IfNoValueF{#5}{%
            \linespread{1.0}\setlength{\parindent}{2em}
            \par\smallskip
            #5
        }
        \vspace{-10pt}
    \end{figure}
}

\crefname{figure}{图}{图}
\crefname{table}{表}{表}
\crefname{equation}{式}{式}
\crefname{chapter}{第\,章}{第\,章}
\crefname{section}{第\,节}{第\,节}
\crefname{subsection}{第\,小节}{第\,小节}
\crefname{subsubsection}{第\,小小节}{第\,小小节}
\crefname{paragraph}{第\,段}{第\,段}
\crefname{appendix}{附录}{附录}
\ifcauc@graduate
\crefname{algocf}{算法}{算法}
\fi
\crefname{subfigure}{子图}{子图}
\crefname{subtable}{子表}{子表}
\Crefname{figure}{图}{图}
\Crefname{table}{表}{表}
\Crefname{equation}{式}{式}
\Crefname{chapter}{第\,章}{第\,章}
\Crefname{section}{第\,节}{第\,节}
\Crefname{subsection}{第\,小节}{第\,小节}
\Crefname{subsubsection}{第\,小小节}{第\,小小节}
\Crefname{paragraph}{第\,段}{第\,段}
\Crefname{appendix}{附录}{附录}
\ifcauc@graduate
\Crefname{algocf}{算法}{算法}
\fi
\Crefname{subfigure}{子图}{子图}
\Crefname{subtable}{子表}{子表}

\ifcauc@graduate
\crefname{nomenclature}{符号说明}{符号说明}
\Crefname{nomenclature}{符号说明}{符号说明}
\fi

\renewcommand{\appendixname}{附录}
\renewcommand{\appendixtocname}{附录}

\g@addto@macro\appendix{%
    \renewcommand{\chaptername}{附录}%
    \ifcauc@bachelor\else
        \renewcommand{\thechapter}{\Alph{chapter}}%
    \fi
}
\ifcauc@graduate
\renewcommand{\listalgorithmcfname}{\centerline{\heiti\zihao{3}算法列表}}
\fi

\newcommand{\makereviewpage}{%
    \ifcauc@graduate
        \newpage
        \thispagestyle{empty}
        \begin{center}
            \vspace*{2cm}
            {\zihao{3}\heiti 答辩委员会评语}\\[2cm]
            \begin{minipage}{0.85\textwidth}
                \zihao{4}\songti
                \vspace{8cm}
            \end{minipage}\\[1cm]
            {\zihao{4}答辩委员会主席签字：\underline{\hspace{4cm}}}\\[1.5cm]
            {\zihao{4}日期：\underline{\hspace{4cm}}}
        \end{center}
    \fi
}

\newcommand{\cauccheck}{%
    \typeout{===== CAUC Thesis Template Configuration =====}%
    \typeout{Degree type: \cauc@degree}%
    \typeout{Bachelor: \ifcauc@bachelor true\else false\fi}%
    \typeout{Graduate: \ifcauc@graduate true\else false\fi}%
    \typeout{Doctor: \ifcauc@doctor true\else false\fi}%
    \typeout{Master: \ifcauc@master true\else false\fi}%
    \typeout{Master-Academic: \ifcauc@masteracademic true\else false\fi}%
    \typeout{Master-Professional: \ifcauc@masterprofessional true\else false\fi}%
    \typeout{Period: \ifcauc@proposal proposal\else final\fi}%
    \typeout{Student Name: \cauc@StudentName}%
    \typeout{Title: \cauc@Title}%
    \typeout{Local Fonts: \ifcauc@localfonts@available true\else false\fi}%
    \typeout{Local Fonts Option: \ifx\cauc@localfonts\@empty auto\else\cauc@localfonts\fi}%
    \typeout{Draft Mode: \ifcauc@draft true\else false\fi}%
    \typeout{Line Numbers: \ifcauc@lineno true\else false\fi}%
    \typeout{Blind Review: \ifcauc@blind true\else false\fi}%
    \typeout{PDF/A: \ifcauc@pdfa true\else false\fi}%
    \typeout{Two-side: \ifcauc@twoside true\else false\fi}%
    \typeout{Open Right: \ifcauc@openright true\else false\fi}%
    \typeout{Engine: \ifxetex XeLaTeX\else\ifluatex LuaLaTeX\else pdfLaTeX\fi\fi}%
    \typeout{Font Config: \ifcauc@localfonts@available local (fonts/ directory)\else system fonts\fi}%
    \typeout{Math Font: \ifx\cauc@mathfont\@empty cm\else\cauc@mathfont\fi}%
    \typeout{Bibliography Style: \ifcauc@bachelor gbt7714-2005-numerical\else gbt7714-numerical\fi}%
    \typeout{Line Spacing: \ifx\cauc@linespacing\@empty 1.389\else\cauc@linespacing\fi}%
    \typeout{Abstract Line Spacing: \cauc@abstractspacing}%
    \typeout{Page Margins: top=\cauc@topmargin, bottom=\cauc@bottommargin, left=\cauc@leftmargin, right=\cauc@rightmargin}%
    \typeout{PDF/A Mode: \ifcauc@pdfa enabled\else disabled\fi}%
    \typeout{Widow/Club Penalty: \the\widowpenalty/\the\clubpenalty}%
    \typeout{Page Alignment: raggedbottom}%
    \typeout{Package: natbib \@ifundefined{natbibversion}{unknown}{\natbibversion}}%
    \typeout{Package: booktabs \@ifundefined{booktabs@version}{unknown}{\booktabs@version}}%
    \typeout{===== End Configuration =====}%
}

\newcommand{\makefrontpages}{%
    \makecover
    \makefrontmatter
}

\@ifundefined{addbibresource}{%
    \let\cauc@orig@bibliography\bibliography
    \renewcommand{\bibliography}[1]{%
        \cauc@orig@bibliography{#1}%
        \addcontentsline{toc}{chapter}{参考文献}%
    }%
}{%
}
